\section*{Chapter \ref{ch:fm:case-studies-ii} --- Case Studies II}

\begin{solution}{exo:fm:case-studies-ii:1}
$97.9/(97.9+266.7)=26.9\%$; $400/266.7=\$1.50$ a share, when book value per share after the loss was still \$10.80.
\end{solution}

\begin{solution}{exo:fm:case-studies-ii:2}
See \cref{def:fm:case-studies-ii:dynamic}. Static initial margin is fixed in dollars at inception: as the position's value rises, the same dollars
are a smaller percentage of it, so the client's leverage with the broker rises exactly as the concentration grows.
\end{solution}

\begin{solution}{exo:fm:case-studies-ii:3}
$10\,000\times(80\,000-48\,078)=\$319$ million; cash runs out at $48\,078+300\,000\,000/10\,000=\$78\,078$.
\end{solution}

\begin{solution}{exo:fm:case-studies-ii:4}
See the proof of \cref{prop:fm:case-studies-ii:zero}. A standstill with a slower joint sale leaves the total fall at most unchanged and spreads it by
agreement; without one, each creditor gains by selling first, which moves its share of the fall onto the others.
\end{solution}

\begin{solution}{exo:fm:case-studies-ii:5}
$0.531\times0.5/7=3.8\%$ for the first, $0.531\times5.5/7=41.8\%$ for the sixth, $0.531\times6.5/7=49.3\%$ for the last; a single broker selling the
whole position suffers an average fall of $0.531/2=26.6\%$.
\end{solution}

\begin{solution}{exo:fm:case-studies-ii:6}
A limit on balances at any one venue, a daily sweep of excess back to a bank or \omterm{def:fm:the-asset-manager-and-the-fund:admin}{custodian}, and a requirement for evidence of segregation. The
exchange's terms of service and financial statements, and any attestation of reserves against liabilities, would have been the records to ask for;
withdrawal delays are the late signal.
\end{solution}

\begin{solution}{exo:fm:case-studies-ii:7}
The first three places lose nothing; the fourth to seventh lose \$1.12, 2.41, 3.70 and 4.99 billion. The last still loses \$4.99 billion.
\end{solution}

\begin{solution}{exo:fm:case-studies-ii:8}
In a race a broker's loss depends on its place and on what the client holds elsewhere: in the calibrated model the last seller needs 49.3\%, and a
20\% margin still loses \$4.99 billion. A margin must scale with concentration against market volume and with what the broker does not see.
\end{solution}

\begin{solution}{pb:fm:case-studies-ii:1}
\begin{enumerate}
\item See \cref{def:fm:case-studies-ii:dynamic}.
\item See \cref{def:fm:case-studies-ii:race}.
\item See \cref{prop:fm:case-studies-ii:zero}.
\item Together they bear a fixed fall, lower if they sell slowly; alone, each can move its share onto the others by selling first.
\item See \cref{dat:fm:case-studies-ii:record}: over \$460 million lost in 45 minutes on 1 August 2012; \$400 million of preferred stock convertible
into about 73\% of the shares sold on 6 August without a shareholder vote; the merger with GETCO on 1 July 2013 at \$3.75 a share or a third of a
new share.
\item \$120 billion of gross exposure and \$9--10 billion of equity, and a request for a standstill; they declined, and sold, some in blocks on 26
March, three under a managed liquidation into April.
\item Assets of 6.6\% of liquid-token payables and 22.6\% of all payables at the petition; customers withdrew a net \$7.0 billion from 1 to 11
November, \$3.3 billion on the 7th, and withdrawals were paused on the 8th.
\item It suspended trading and cancelled all trades since midnight on 8 March; the courts held that lawful, and the FCA fined it
\pounds9\,245\,900 for failing to maintain orderly trading.
\item \Cref{tab:fm:case-studies-ii:rescue}: old shareholders from 100\% to 26.9\%, book value per share from \$15.29 to \$10.80 to \$4.00.
\item Seven brokers of \$17 billion each, linear permanent impact, the coefficient 0.531 chosen so that the sixth with 9.4\% loses \$5.5 billion.
\item \$0, 0.34, 1.63, 2.92, 4.21, 5.50 and 6.79 billion.
\item \$0, 0, 0, 1.12, 2.41, 3.70 and 4.99 billion; 3.8\%, 11.4\%, 19.0\%, 26.6\%, 34.2\%, 41.8\% and 49.3\%.
\item Unequal exposures, the price fall before default, short positions and hedges, block trades that do not move the screen price linearly, and
recovery from the client's estate.
\item \$319 million at \$80\,000, \$533 million at the peak; cash of \$300 million runs out at \$78\,078.
\item Dynamic margins with concentration and liquidity add-ons, recomputed daily, and the right to raise them quickly.
\item Ask the client, contractually, for its positions across brokers; watch public filings; assume the worst when it will not say.
\item A limit per venue and a daily sweep of excess, with tighter limits where segregation cannot be verified.
\item Hold capital for a morning like 1 August 2012 under a hard aggregate exposure limit, and arrange committed funding that can be drawn in a day.
\item At 9.4\% margin, \$0 for the first seller up to \$6.79 billion for the last (\$4.99 billion at 20\%); the old shareholders kept 26.9\% of the
2012 firm, \$391 million of book equity instead of \$1\,057 million.
\item Margin concentrated clients dynamically, for the whole position's exit rather than one broker's share, and agree in advance how to act if the
client fails. Cap the firm's own balances at any venue and hold capital and committed funding for a loss of a day's worth of its largest exposure.
\end{enumerate}
\end{solution}

\begin{solution}{iq:fm:case-studies-ii:1}
Static margin is fixed at inception as a percentage of notional; dynamic margin is recomputed from the portfolio's current volatility,
concentration, liquidity and bias.
\iqlookfor{erosion of static margin as prices move.}
\end{solution}

\begin{solution}{iq:fm:case-studies-ii:2}
Pre-trade limits, an aggregate exposure limit with a kill switch, and deployment checks; what it costs the owners is set by the capital and funding
held against such a loss and whether a rescue must be arranged in days.
\iqlookfor{controls, then capital.}
\end{solution}

\begin{solution}{iq:fm:case-studies-ii:3}
Weigh the race: selling first moves the loss to others, but a joint slow sale lowers the total; join a managed liquidation if it is binding and
covers the overlapping positions, and hedge meanwhile.
\iqlookfor{the zero-sum structure of the race.}
\end{solution}

\begin{solution}{iq:fm:case-studies-ii:4}
Only what trading needs that day, swept back daily, under a limit set by what the firm could lose without distress, since the balance is an
unsecured claim unless segregation is verified.
\iqlookfor{the balance as a credit exposure.}
\end{solution}

\begin{solution}{iq:fm:case-studies-ii:5}
That the exchange may cancel the trades under its rules, as the nickel market's did; book nothing until the trades are confirmed to stand.
\iqlookfor{cancellation powers.}
\end{solution}

\begin{solution}{iq:fm:case-studies-ii:6}
With linear permanent impact the total fall is fixed at half the impact coefficient times the position; the order decides who bears it, and the
last seller bears the most. Slower joint selling or block trades can lower the total.
\iqlookfor{the integral argument.}
\end{solution}
